% Generated draft: requires mathematical and source review.
\documentclass[11pt]{article}
\usepackage[margin=25mm]{geometry}
\usepackage{fontspec}
\setmainfont{Latin Modern Roman}
\usepackage{amsmath,amssymb,mathtools,unicode-math}
\setmathfont{Latin Modern Math}
\usepackage{enumitem,needspace,xurl,hyperref}
\setlength{\parindent}{0pt}
\setlength{\parskip}{5pt}
\setlength{\emergencystretch}{3em}
\newcommand{\N}{\mathbb{N}}
\newcommand{\Z}{\mathbb{Z}}
\newcommand{\Q}{\mathbb{Q}}
\newcommand{\R}{\mathbb{R}}
\newcommand{\C}{\mathbb{C}}
\newcommand{\F}{\mathbb{F}}
\newcommand{\A}{\mathbb{A}}
\newcommand{\PP}{\mathbb{P}}
\newcommand{\E}{\mathbb{E}}
\newcommand{\eps}{\varepsilon}
\newcommand{\dd}{\,\mathrm{d}}
\begin{document}
\section*{UnsolvedMath ID: 20001860}
\textbf{Source problem code:} catalogue entry
\subsection{Definitions and mathematical statement}
Owen's source asks when the Galois group of a particular rigid framework is a subgroup of the generic graph Galois group, written \(G(G,p)\subseteq G(G)\), including the proposed isostatic case. It defines generic coordinates as algebraically independent over \(\Q\), but does not define the two Galois groups or the normalization. The parenthetical ``Isostatic is sufficient but not necessary'' is retained as source wording, not assumed as a verified universal theorem.

For the aid's full-fiber interpretation, take a labelled planar Laman graph, fix Euclidean gauge algebraically, and use its squared-distance equations to obtain a generically finite separable realization system over the field \(K\) of generic squared edge labels. Here Laman means \(|E|=2|V|-3\) and each vertex subset of size at least two spans at most \(2|V'|-3\) edges. Let \(L/K\) be the splitting field of the full normalized generic fiber, and \(\mathcal G=\operatorname{Gal}(L/K)\). This choice concerns all normalized realizations, not only the coordinate field of one selected framework.

The reported good-specialization criterion is: specialize to a label value \(b\) away from poles, degree loss, and the branch locus of a suitable finite model. Then the Galois group of the specialized full realization fiber over its label field is isomorphic to a decomposition subgroup of \(\mathcal G\), hence embeds in \(\mathcal G\). Choosing a place over \(b\) determines the subgroup only up to conjugacy. In the associated Galois normalization, write \(D\leq\mathcal G\) for the decomposition group and \(I\triangleleft D\) for inertia. The canonical residue-field group is
\[
 \operatorname{Gal}(\kappa(\widetilde b)/\kappa(b))\cong D/I.
\]
At an unramified good specialization \(I=1\); at a branch point the canonical comparison supplies a subquotient, not a distinguished subgroup embedding. Further identification with a chosen specialized realization-field convention must be checked. A nonsingular selected realization does not itself show that the entire label value is outside the branch locus.

The aid supplies the following local/global counterexample to that certification step. Take the five-vertex type-1 Laman graph
\[
 E=\{12,13,23,14,24,35,45\},
\]
with rational squared edge labels
\[
 \ell_{12}^2=\ell_{13}^2=\ell_{24}^2=1,\qquad
 \ell_{23}^2=\ell_{14}^2=2,\qquad
 \ell_{35}^2=\ell_{45}^2=5/4.
\]
In its normalized fiber there is an isostatic realization with nonzero squared-distance Jacobian determinant, reported as \(64\), and another ramified realization with determinant \(0\), with no coincident vertices or zero bars. The value \(64\) depends on equation and coordinate normalization; the invariant distinction is full rank versus rank loss. Isostatic here means an infinitesimally rigid realization with the minimal independent bar count after removing Euclidean motions. The example proves that isostaticity at one selected point cannot replace the global good-fiber hypothesis. It does not by itself prove that the specialized full-fiber Galois group fails to embed abstractly, nor settle every intended definition of Owen's question.
\subsection{Short English statement}
Away from poles, degree loss and branching, the full realization-fiber Galois group specializes to a decomposition subgroup of the generic group. At branching the canonical comparison is a quotient by inertia. The aid's Laman example has both a nonsingular and a ramified realization over the same labels, so one isostatic framework does not certify a good entire fiber.
\subsection{Sources}
\begin{itemize}
\item Local UnsolvedMath registry, permanent ID \(20001860\), source code \texttt{catalogue entry}; canonical collection \texttt{aim-geometry-notes.json}, notes[197].
\item ``Configuration Spaces of Linkages: Open Problems,'' notes by Elissa Ross, AIM workshop, 25--31 October 2014, problem \(11\), attributed to Owen: \url{https://aimath.org/pastworkshops/linkagesproblems.pdf}. The supplied verification corrects the extraction's problem ``1'' and extraneous final ``1''; the permanent ID and source code remain unchanged. The PDF was not opened in this drafting session.
\item Supplied research attempt \(1\), \texttt{partial\_result}, run status ``open,'' rigor label ``proved,'' novelty confidence ``low.'' Registry status \texttt{partially\_solved}, classification \texttt{PARTIAL-PROGRESS}, updated 14 August 2026. Original remarks and literature are empty; no full authoritative specialization reference is supplied, and proposer/year registry fields are null despite the Owen attribution.
\end{itemize}
\end{document}
